\answer{ex:01:1}{(a)~$8\cdot10^{10}$ \nt{GLUT}: $10$ populations de la Terre; ceux de diffusion, $\approx1.9\cdot10^6$: une ville de la taille de Vienne. (b)~$\approx8\cdot10^{14}$ terminaisons, dont $\approx3\cdot10^{11}$ de diffusion, soit une part de $\approx4\cdot10^{-4}$, $\approx16$ fois plus grande que la part de ces neurones ($2.2\cdot10^{-5}$).}
\answer{ex:01:2}{(a)~$\tau_c\ln(c_0/c_*)\approx4.6\ms$. (b)~La ligne est libre si $T\ge\tau_c\ln101$: jusqu'à $\approx22$ libérations par seconde au lieu de $\approx220$.}
\answer{ex:01:3}{$V(s_k)=\gamma^{K-1-k}r$; $\delta_{\text{lampe}}=\gamma^Kr=re^{-T/\tau_\gamma}\approx0.60\,r$ quel que soit $\Delta$, $\tau_\gamma\approx1.95\,\text{s}$.}
\answer{ex:01:4}{$n^*=93$, $47$, $25$, $12$: $n^*\approx5/\alpha$ pour $\alpha$ petit.}
\answer{ex:01:5}{\nt{GLUT}: $1\to10\to10^4\to10^7$, soit $\approx10^7$ neurones, $4$ maillons, $4\ms$. \nt{DOPA}: $1\to10\to3.2\cdot10^5$, soit $\approx3\cdot10^5$ neurones, $3$ maillons, $30$ fois moins; c'est le même ordre de grandeur que le nombre de neurones \nt{DOPA}.}
\answer{ex:01:6}{$\lambda=f_EJ_E-f_IJ_I$, d'où $J_I=37.5$; rapport des courants $f_IJ_I/f_EJ_E=0.94$; avalanche ($\lambda\ge1$) dès que l'inhibition est affaiblie de seulement $6.7\,\%$.}
\answer{ex:01:7}{$n\ge57$ présentations pour chaque délai dans chaque condition. La même décroissance pourrait venir, par exemple, d'une imprécision de l'estimation interne du temps qui croît avec $T$.}
